\section{Alternativas e limites}

\begin{frame}{Software livre e padrões abertos podem ampliar opções}
  \begin{block}{Possíveis contribuições}
    \begin{itemize}
      \item Estudo, auditoria e modificação do software.
      \item Redistribuição e adaptação conforme a licença \parencite{stallman2002free}.
      \item Interoperabilidade e maior portabilidade, quando efetivamente implementadas.
    \end{itemize}
  \end{block}
  \begin{alertblock}{O que não garantem isoladamente}
    Capacidade operacional, independência de hardware, ausência de custos, segurança automática ou governança adequada.
  \end{alertblock}
  \note{Use fontes primárias para definir software livre e fontes técnicas para discutir padrões. Evite apresentar software livre como solução suficiente para todos os problemas de soberania.}
\end{frame}
